\documentclass[]{article}

%\usepackage[authoryear]{natbib}
%\usepackage[round,sort&compress,sectionbib]{natbib}

\usepackage[style=authoryear,maxnames=5,backend=bibtex,natbib=true]{biblatex}
%\bibliographystyle{plain}
\addbibresource{AIP-GalClusTurb.bib}

%\bibliographystyle{plainnat}
\usepackage{hyperref}
\usepackage{siunitx}
\usepackage{bm}
\let\vec\bm
\usepackage{xurl}
\usepackage{subcaption}
\usepackage{graphicx}   % Including figure files
\usepackage{amsmath}   % Advanced maths commands
\usepackage{amssymb}   % Extra maths symbols
\usepackage{xcolor}
\usepackage[makeroom]{cancel}
%\usepackage{soul}
\newcommand{\mathcolorbox}[2]{\colorbox{#1}{$\displaystyle #2$}}
\colorlet{lgray}{gray!50!white}
\colorlet{lred}{red!50!white}
\newcommand{\bnabla}{\bm{\nabla}}
\newcommand{\bcdot}{\bm{\cdot}}
\newcommand{\ud}{\,\mathrm{d}}


%opening
\title{Turbulent energy/pressure ratio}
\author{Lorenzo M. Perrone, TB, EP, CP}

\begin{document}

\maketitle

\begin{abstract}
In this document we define the turbulent ratio.
\end{abstract}


\section{Definitions}

\subsection{Thermodynamic quantities}

The volumetric internal energy density $\epsilon_{\mathrm{int}}$ (in erg cm$^{-3}$) is defined as:
%
\begin{align}
	\epsilon_{\mathrm{int}} = \frac{p}{\gamma - 1},
\end{align}
%
where $p$ is the thermal pressure (erg cm$^{-3}$) and $\gamma$ is the adiabatic index. The mass density $\rho$ (in g cm$^{-3}$) is defined as,
%
\begin{align}
	\rho = n \mu m_\mathrm{u} ,
\end{align}
%
with $n$ the particle number density (in cm$^{-3}$), $\mu$ the mean molecular weight, and $m_\mathrm{u}$ the atomic mass unit. The thermal pressure $p$ is related to the temperature $T$ and the number density $n$ by the ideal gas law:
%
\begin{align}
	p = n k_\mathrm{B} T = \frac{\rho k_{\mathrm{B}}T}{m_\mathrm{u} \mu}.
\end{align}
%
%\subsection{Sound speed}
%
The adiabatic sound speed $c_{\mathrm{s}}$ is defined as
%
\begin{align}\label{eq:sound_speed}
	c_{\mathrm{s}}^2 \equiv \frac{\gamma p}{\rho} = \gamma (\gamma - 1) \frac{\epsilon_{\mathrm{int}}}{\rho} = \frac{\gamma k_{\mathrm{B}}T}{m_\mathrm{u} \mu}.
\end{align}
%

\subsection{Volume averages}

We define volume averages $\langle p \rangle_V$ or $\overline{p}$ for convenience:
%
\begin{align}
	\langle p \rangle_V = \overline{p} = \frac{1}{V} \int_V p \ud^3 x.
\end{align}

\subsection{Energy densities and energies}

All the quantities have the same meaning as in the main paper, i.e. $\varepsilon_{\mathrm{k,turb}}$ is the turbulent kinetic energy density, and $\mathcal{E}_{\mathrm{k,turb}} = \int_V \varepsilon_{\mathrm{k,turb}} \ud^3 x$ is its volume integral. In terms of the volume average:
%
\begin{align}
	\mathcal{E}_{\mathrm{k,turb}} = V \overline{\varepsilon}_{\mathrm{k,turb}}.
\end{align}
%
The thermal energy $E_{\mathrm{th}}$ in a volume $V$ is:
%
\begin{align}
	E_{\mathrm{th}} =  \int_V \epsilon_{\mathrm{int}} \ud^3 x = \frac{1}{\gamma - 1} V \overline{p}.
\end{align}
%
%where $\ud m = \rho \ud^3 x$.

\subsection{Turbulent pressure}

We define the turbulent kinetic pressure $p_{\mathrm{turb}}$ in the filtering approach as 
%
\begin{align}\label{eq:turb-press}
	p_{\mathrm{turb}} = \frac{2}{3} \varepsilon_{\mathrm{k,turb}},
\end{align}
%
while its mean value over a volume $V$ is given by:
%
\begin{align}
	\overline{p}_{\mathrm{turb}} \equiv \frac{2}{3} \frac{\mathcal{E}_{\mathrm{k,turb}}}{V} = \frac{2}{3} \overline{\varepsilon}_{\mathrm{k,turb}}.
\end{align}



\subsection{Velocity dispersion}

We introduce the local three-dimensional velocity dispersion $\sigma_{\mathrm{3D}} (\bm x)$ as:
%
\begin{align}\label{eq:3d-vel-disp}
	\sigma_{\mathrm{3D}}^2 (\bm x)= \sum_i \mu_2(u_i, u_i),
\end{align}
%
where $i=x,y,z$, and the dependence on the filter size $\ell$ is implicit in the definition of $\mu_2$. This is a local quantity that depends on $\bm x$. Its average over $V$ gives an rms value of $\sigma_{\mathrm{3D},\mathrm{rms}}$:
%
\begin{align}
	\sigma_{\mathrm{3D},\mathrm{rms}}^2  = \frac{1}{V} \int_V \sigma_{\mathrm{3D}}^2 (\bm x) \ud^3 x = \overline{\sigma_{\mathrm{3D}}^2}.
\end{align}
%
We define the one-dimensional rms velocity dispersion
%For isotropic turbulence we define the one-dimensional rms velocity dispersion
$\sigma_{\mathrm{1D},\mathrm{rms}}$ as
%
\begin{align}
	\sigma_{\mathrm{1D},\mathrm{rms}}^2 \equiv \overline{\mu_2(u_x, u_x)} = \overline{\mu_2(u_y, u_y)} = \overline{\mu_2(u_z, u_z)},
\end{align}
%
which for isotropic turbulence is related to the three-dimensional rms velocity dispersion by:
%
\begin{align}
		\sigma_{\mathrm{3D},\mathrm{rms}}^2 = 3 \sigma_{\mathrm{1D},\mathrm{rms}}^2.
\end{align}




\section{Homogeneous subsonic turbulence}

We assume that the ICM is at rest and in hydrostatic equilibrium, with background hydrodynamic profiles $\rho_0 (\bm x)$, $p_0 (\bm x)$, $T_0 (\bm x)$, and small fluctuations $\bm u (\bm x)$, $\delta \rho (\bm x), \ldots$ on top of it.

\subsection{Subsonic turbulence}

If the turbulence is subsonic ($\mathcal{M} =  u / c_{\mathrm{s}} \ll 1$), then it can be shown via standard derivations that 
%
\begin{align}
	\frac{\delta \rho}{\rho_0} \sim \mathcal{O} \left(\frac{u}{c_{\mathrm{s}}}\right) \sim \mathcal{M}.
\end{align}
%
In this case, after filtering $\delta \rho_\ell / \langle \rho \rangle_\ell \sim \mathcal{O} (\mathcal{M})$, and the turbulent kinetic energy density reduces to\footnote{Will be shown in an Appendix in the paper.}
%
\begin{align}
	\varepsilon_{\mathrm{k,turb}} &\equiv \frac{1}{2} \langle \rho \rangle_\ell \sum_i \mu_2(u_i, u_i) + \frac{1}{2} \sum_i \mu_3(\rho, u_i, u_i) \nonumber \\ 
	&\simeq \frac{1}{2} \langle \rho_0 (\bm x) \rangle_\ell \sigma_{\mathrm{3D}}^2 (\bm x), 
\end{align}
%
where we used the definition of the three-dimensional velocity dispersion Eq.~\eqref{eq:3d-vel-disp}.

\subsection{Homogeneous turbulence}

We can further assume that the turbulence is homogeneous in a certain region of interest. This means that the background quantities and the statistical properties of the fields do not depend on $\bm x$. This implies that $\langle \rho \rangle_\ell \simeq \rho_0 = \text{const}$ and the turbulent kinetic energy density reduces to
%
\begin{align}
	\varepsilon_{\mathrm{k,turb}} (\bm x) = \frac{1}{2} \rho_0 \sigma_{\mathrm{3D}}^2 (\bm x).
\end{align}
%


\subsection{Turbulent pressure}

From the definition of the turbulent pressure Eq.~\ref{eq:turb-press} it follows that
%
\begin{align}
	\overline{p}_{\mathrm{turb}} \hspace{4em}  &= \frac{2}{3} \overline{\varepsilon}_{\mathrm{k,turb}} \nonumber \\
	\begin{tabular}{p{6em}}
		subsonic \\ turbulence
	\end{tabular}    &\simeq \frac{\cancel{2}}{3} \frac{1}{\cancel{2}} \overline{\langle \rho \rangle_\ell \sigma_{\mathrm{3D}}^2 (\bm x)} \nonumber \\
	\begin{tabular}{p{6em}}
		homogeneous \\ turbulence
	\end{tabular}    &\simeq \frac{1}{3}  \rho_0 \sigma_{\mathrm{3D,rms}}^2 \nonumber \\
	\begin{tabular}{p{6em}}
	isotropic \\ turbulence
	\end{tabular}    &\simeq   \rho_0 \sigma_{\mathrm{1D,rms}}^2 
\end{align}


\section{Turbulent ratios}

\subsection{Definition based on filtered turbulent energy}

We define the mean turbulent pressure/energy ratio $f_{\mathrm{turb}}$ in a volume $V$ as:
%
\begin{align}
	f_{\mathrm{turb}} &\equiv \frac{\mathcal{E}_{\mathrm{k,turb}}}{\mathcal{E}_{\mathrm{k,turb}} + E_{\mathrm{th}}} =  \frac{\frac{3}{2} \cancel{V} \overline{p}_{\mathrm{turb}}}{\frac{3}{2} \cancel{V} \overline{p}_{\mathrm{turb}} + \frac{1}{\gamma - 1} \cancel{V} \overline{p}_{\mathrm{th}}} \nonumber \\
	&= \frac{ \overline{p}_{\mathrm{turb}}}{ \overline{p}_{\mathrm{turb}} + \frac{2/3}{\gamma - 1}  \overline{p}_{\mathrm{th}}} \xrightarrow{\gamma = 5/3} \frac{ \overline{p}_{\mathrm{turb}}}{ \overline{p}_{\mathrm{turb}} +   \overline{p}_{\mathrm{th}}}.
\end{align}
%
We have ignored magnetic pressure and cosmic ray pressure.

\subsection{Comparison with existing formulae in literature}

To compare with existing formulae in the literature, we take the homogeneous subsonic limit. This allows us to take the density out of the volume average, and use rms values for the velocity dispersion. Implicitly, this is what XRISM and observational estimates do (i.e., take a mean density, and multiply by the mean velocity dispersion, ignoring correlations between the local values).
%
For an adiabatic gas, and subsonic, homogeneous turbulence, the turbulent ratio can be rewritten as:
%
\begin{align}
	f_{\mathrm{turb}} &= \frac{ \overline{p}_{\mathrm{turb}}}{ \overline{p}_{\mathrm{turb}} +   \overline{p}_{\mathrm{th}}} \simeq \frac{\frac{1}{3}  \rho_0 \sigma_{\mathrm{3D,rms}}^2}{\frac{1}{3}  \rho_0 \sigma_{\mathrm{3D,rms}}^2 + \frac{\rho_0 k_{\mathrm{B}}\overline{T}}{m_\mathrm{u} \mu}} \nonumber \\
	&= \frac{\sigma_{\mathrm{3D,rms}}^2}{\sigma_{\mathrm{3D,rms}}^2 + \frac{3  k_{\mathrm{B}}\overline{T}}{m_\mathrm{u} \mu}} \begin{tabular}{p{12em}}
		\citep[][Eq.10]{Eckert2019b}
	\end{tabular}
\end{align}
%
Using the expression in Eq.~\ref{eq:sound_speed} for the mean adiabatic sound speed
%
\begin{align}
	\overline{c}_\mathrm{s}^2 = \frac{\gamma k_\mathrm{B} \overline{T}}{\mu m_u},
\end{align}
%
and defining the three-dimensional mean turbulent mach number $\overline{\mathcal{M}}_{\mathrm{3D}}^2 = \sigma_{\mathrm{3D,rms}}^2 / \overline{c}_\mathrm{s}^2$, the expression above can alternatively be rewritten as:
%
\begin{align}
	f_{\mathrm{turb}} \simeq \frac{\overline{\mathcal{M}}_{\mathrm{3D}}^2}{\overline{\mathcal{M}}_{\mathrm{3D}}^2 + \frac{3}{\gamma}}  \begin{tabular}{p{20em}}
		\citep[][Eq.1]{XRISMCollaboration2025a} and \citep[][Eq.1]{XRISMCollaboration2025c}
	\end{tabular}.
\end{align}
%
If the turbulence is isotropic, the turbulent ratio can be expressed in terms of the one-dimensional mean turbulent mach number $\overline{\mathcal{M}}_{\mathrm{1D}}^2 = \sigma_{\mathrm{1D,rms}}^2 / \overline{c}_\mathrm{s}^2$ as:
%
\begin{align}
	f_{\mathrm{turb}} \simeq \frac{\overline{\mathcal{M}}_{\mathrm{1D}}^2}{\overline{\mathcal{M}}_{\mathrm{1D}}^2 + \frac{1}{\gamma}}  \begin{tabular}{p{20em}}
		\citep[][Eq.1]{Zhang2025a} 
	\end{tabular}.
\end{align}

\printbibliography

\end{document}
